\section{Related Work} \label{section:related_work}

\textbf{Current concentration inequalities for the variance.} Upper and lower inequalities for the variance were presented in \citet[Theorem 10]{maurer2009empirical}. They are based on the concentration of self-bounding random variables \citep{maurer2006concentration}. Another concentration inequality for the variance can be found in the proof of~\citet[Theorem 1]{audibert2009exploration}, which decouples the analysis into those of the mean and second centered moment, in a similar spirit to our contribution. However, these inequalities rely on conservatively upper bounding the \emph{variance of the empirical variance}, thus being loose (this is also the case  for \citet[Theorem 12]{voravcek2025star}).  In contrast, our inequalities empirically estimate the variance of the empirical variance, resulting in tighter confidence sets.  We defer an extended comparison of these inequalities to Section~\ref{section:comparison}. Less closely related to our work, other inequalities rely on the Kolmogorov-Smirnov distance between the empirical distribution and a reference distribution. % Within this line of work, \citet{austern2022efficient} had originally adapted the finite-sample results from \citet{romano2000finite} to obtain variance bounds that are often significantly tighter than those from \citet{maurer2009empirical}; however, these were outperformed and replaced by the inequalities introduced in this contribution. 


% Less closely related to our work, other inequalities rely on the Kolmogorov-Smirnov (KS) distance between the empirical distribution and a second distribution. For instance,~\citet[Lemma 2]{austern2022efficient} leveraged the concentration inequalities from \citet{romano2000finite}, which build on the KS distance between the empirical distribution and the actual distribution, in order to derive inequalities for the variance.
% (a similar approach could probably be derived by means of Berry-Esseen bounds, i.e. understanding the KS distance of the normalized empirical distribution and the standard normal distribution). \citet[Appendix D]{austern2022efficient} elucidated the use of these inequalities for variance estimation in the case of independent and identically distributed (iid) random variables. However, these methods strongly rely on independence assumptions and are also tailored to the batch setting, where the sample size is fixed in advance.\footnote{In particular, the ``empirical Berry-Esseen'' inequalities from \citet[Appendix D]{austern2022efficient} combine the inequalities from \citet{romano2000finite} and \citet{maurer2009empirical}. Our contribution improves on the latter, hence opening the door for also improving the results in \citet{austern2022efficient}.}
% making them unsuitable for the online setting (where iid is generally too strong of an assumption, and practitioners seek anytime valid guarantees that allow them to conduct experiments sequentially).

 

\textbf{Empirical Bernstein inequalities for the mean.} Based on combining Bennett's inequality and upper concentration inequalities for the variance, 
\citet[Theorem 11]{maurer2009empirical} proposed a well known empirical Bernstein inequality for the mean, improving a similar inequality presented in \citet[Theorem 1]{audibert2009exploration}. These inequalities are not asymptotically sharp as presented (in that the first order limiting width does not match that of the oracle Bernstein inequality, including constants). However, they can be amended to recover sharpness if the probability split of the union bound is carefully designed (as pointed out recently in \citet[Section B.2]{wang2024sharp}). Nevertheless, their bounds were empirically significantly looser \citep[Figure 3]{waudby2024estimating} than those presented in \citet[Theorem 4]{howard2021time}  and \citet[Theorem 2]{waudby2024estimating}, which are also sharp. Related contributions include \citet[Theorem 2]{mnih2008empirical}, \citet[Corollary 4]{jang2023tighter}, \citet[Theorem 3]{orabona2023tight}, and \citet[Corollary 1]{martinez2024empirical}. 


% \paragraph{Betting-based concentration.} A profound relationship exists between concentration inequalities and the regret guarantees of online learning algorithms with linear losses, as established by \cite{rakhlin2017equivalence}. Consequently, the study of concentration phenomena through the lens of gambling-based frameworks has become an active and increasingly influential research direction \citep{jun2019parameter, waudby2024estimating, voravcek2025star}. Empirical Bernstein bounds have also been developed under this line of research, see e.g. \citet[Corollary 4]{jang2023tighter} or \citet[Theorem 3]{orabona2023tight}. Nonetheless, the resulting concentration inequalities are often not available in closed form, and inverting these bounds tends to be analytically opaque and computationally challenging. For this reason, we focus on non-betting closed-form inequalities in this work. However, we highlight that the recent betting-based contribution \citet{voravcek2025star} proposed a novel concentration inequality for the second moment centered at any arbitrary point; we provide additional commentary in Appendix~\ref{section:comparison}.


\textbf{Time-uniform Chernoff inequalities.} Our work falls under the time-uniform Chernoff inequalities umbrella from \citet{howard2020time,howard2021time,waudby2024estimating}. A key proof technique of this line of work is the derivation of sophisticated nonnegative supermartingales, followed by an application of Ville's inequality \citep{ville1939etude}, an anytime-valid version of Markov's inequality. 


